\startcomponent ch03
\product fduthesis

\chapter{公式、图表与代码}

\section{数学公式}

\subsection{行内公式}

行内公式使用 \type{$...$} 标记。
例如，梯度下降的参数更新公式为 $\theta_{t+1} = \theta_t - \alpha \nabla \mathcal{L}(\theta_t)$，
其中 $\alpha$ 为学习率，$\nabla \mathcal{L}$ 为损失函数的梯度。

\subsection{行间公式}

行间公式使用 \type{\startformula} 环境，
通过 \type{\placeformula} 自动编号
（格式为"章号.公式号"）。

交叉熵损失函数定义如下：

\placeformula[eq:crossentropy]
\startformula
  \mathcal{L} = -\frac{1}{N}\sum_{i=1}^{N}\left[y_i\log\hat{y}_i + (1-y_i)\log(1-\hat{y}_i)\right]
\stopformula

其中 $N$ 为样本数量，$y_i \in \{0, 1\}$ 为真实标签，$\hat{y}_i$ 为模型预测概率。

Softmax 函数的定义：

\placeformula[eq:softmax]
\startformula
  \text{softmax}(z_i) = \frac{e^{z_i}}{\sum_{j=1}^{K} e^{z_j}}, \quad i = 1, 2, \ldots, K
\stopformula

多行公式可以使用 \type{\startalign} 环境对齐：

\placeformula
\startformula
  \startalign
    \NC f(x) \NC = (x+1)^2                    \NR
    \NC      \NC = x^2 + 2x + 1               \NR
  \stopalign
\stopformula

\section{表格}

表格标题置于表格上方，编号格式为"表 章号-序号"。
以下是不同模型的性能对比表：

\placetable
  [here][tab:comparison]
  {不同模型的性能对比}
  {\starttabulate[|c|c|c|c|c|]
    \HL
    \NC {\bf 模型} \NC {\bf 准确率} \NC {\bf 精确率} \NC {\bf 召回率} \NC {\bf F1 值} \NC \NR
    \HL
    \NC TextCNN  \NC 0.89 \NC 0.87 \NC 0.85 \NC 0.86 \NC \NR
    \NC BiLSTM   \NC 0.91 \NC 0.90 \NC 0.88 \NC 0.89 \NC \NR
    \NC BERT     \NC 0.94 \NC 0.93 \NC 0.92 \NC 0.92 \NC \NR
    \NC 本文方法 \NC {\bf 0.96} \NC {\bf 0.95} \NC {\bf 0.94} \NC {\bf 0.94} \NC \NR
    \HL
  \stoptabulate}

可以通过 \type{\in{表}{~}[tab:comparison]} 引用表格：
如\in{表}{~}[tab:comparison]所示，
本文方法在各项指标上均优于基线模型。

\section{图片}

图片标题置于图片下方，编号格式为"图 章号-序号"。
将图片文件放入 \type{figures/} 目录，使用以下方式插入：

\starttyping
\placefigure
  [here][fig:architecture]
  {模型整体架构图}
  {\externalfigure[architecture][width=0.8\textwidth]}
\stoptyping

引用方式：\type{\in{图}{~}[fig:architecture]}。

\section{代码}

使用 \type{\starttyping} 环境排版代码块：

\starttyping
import torch
import torch.nn as nn

class TextClassifier(nn.Module):
    def __init__(self, vocab_size, embed_dim, num_classes):
        super().__init__()
        self.embedding = nn.Embedding(vocab_size, embed_dim)
        self.fc = nn.Linear(embed_dim, num_classes)

    def forward(self, x):
        x = self.embedding(x).mean(dim=1)
        return self.fc(x)
\stoptyping

\stopcomponent
